% Do not change document class, margins, fonts, etc.
\documentclass[a4paper,oneside]{scrbook}

% some useful packages
\usepackage[utf8]{inputenc}
\usepackage{graphicx}
\usepackage{latexsym}
\usepackage{amsmath}
\usepackage{amssymb}
\usepackage{tabularx}
\usepackage{booktabs}
\usepackage{algorithm}
\counterwithin{algorithm}{chapter}
\usepackage{algorithmic}
\usepackage{csquotes}
\renewcommand{\algorithmiccomment}[1]{\hfill\textit{// #1}}
\usepackage[usenames,dvipsnames,table]{xcolor}
\usepackage{array}
\usepackage{xurl}
\usepackage[colorlinks,citecolor=Green]{hyperref}

% Table formatting helpers
\definecolor{Missing}{RGB}{192,0,0}
\newcommand{\miss}[1]{#1\textsuperscript{\textcolor{Missing}{**}}}

% example enviroments
\newtheorem{definition}{Definition} \newtheorem{proposition}{Proposition}

\begin{document}

% Quick fix to remove the 0 in front of sections.
\makeatletter
\renewcommand\thesection{\arabic{section}}
\renewcommand\thesubsection{\arabic{section}.\arabic{subsection}}
\renewcommand\thesubsubsection{\arabic{section}.\arabic{subsection}.\arabic{subsubsection}}
\makeatother

\subject{Web Data Integration \\ Project Outline} 
\title{Pharmaceutical Drug Integration: Linking US and EU Approval Data}
\author{
  Gabriel Himmelein (1649181)\\
  Özlem Ölçer (2265825)\\
  Anatole Lobenko (2187322)\\
  Hai Nguyen (2267317)\\ 
  Thu Huyen Trang Trinh (2111810)
} \date{October 6, 2026}
\publishers{{\small Submitted to}\\
  Data and Web Science Group\\
  Dr. Ralph Peters\\
  Prof. Dr. Christian Bizer\\
  University of Mannheim\\}
\maketitle

% ============================================================
\section{Use Case}
% ============================================================

This project aims to build a unified view of approved pharmaceutical drugs in the United
States and the European Union, enabling researchers and analysts to query regulatory
status, approval history, and product attributes across both markets from a single integrated data set. To achieve this, we combine data from three complementary sources.
Drugs@FDA is the FDA's official database of US-approved drugs. EMA Medicines is
the European Medicines Agency's register of EU-authorised medicines. The openFDA
NDC Directory is a supplementary US drug listing with broader product coverage. The
central entity is the active substance, the pharmacologically active ingredient in a drug
product. Active substances serve as a natural integration key because they persist across
different brand names, manufacturers, regulatory jurisdictions, and formulations. The
same substance can even carry different names across regions, as with paracetamol in
the EU and acetaminophen in the US. Beyond regional synonyms, biosimilars carry
manufacturer-specific suffixes that further complicate direct matching. Resolving these
naming variations through identity resolution is therefore a central task of the project,
ultimately producing one clean record per active substance with consolidated attributes
from all three sources.

% ============================================================
\section{Requirements Suitability}
% ============================================================

\subsection{Schema and basic profile}

\begin{table}[H]
\centering
\caption{Overview of the input datasets}
\label{tab:datasets}
\scriptsize
\begin{tabularx}{\textwidth}{|l|l|l|l|l|>{\raggedright\arraybackslash}X|}
\hline
\textbf{Dataset} & \textbf{Source (*)} & \textbf{Format} & \textbf{\# Ent.} & \textbf{\# Attr.} & \textbf{List of attributes (**)} \\
\hline
Drugs@FDA & U.S. FDA & TXT/TSV & 51,864 & 56 & ApplNo, ProductNo, Form, Strength, ReferenceDrug, DrugName, ActiveIngredient, \miss{ApplPublicNotes}, SponsorName, SubmissionType, SubmissionStatus, \miss{SubmissionsPublicNotes}, \miss{ReviewPriority}, MarketingStatusID, TECode, \dots \\
\hline
EMA Medicines & EMA & XLSX & 2,746 & 39 & Name of medicine, EMA product number, Medicine status, \miss{Opinion status}, Active substance, \miss{Species (vet.)}, ATC code, Biosimilar, Generic, Orphan medicine, \miss{Opinion adopted date}, \miss{Revision number}, \dots \\
\hline
openFDA NDC & U.S. FDA & JSON$\rightarrow$CSV & 138,209 & 10 & product NDC, brand name, generic name, active ingredients, manufacturer, dosage form, route, marketing start date, marketing category, application number \\
\hline
\end{tabularx}
\smallskip
\raggedright
{\scriptsize
\textbf{(*) Source:} Downloaded directly from the respective official provider websites. \\
\textbf{(**) Missingness:} \textsuperscript{\textcolor{Missing}{**}} indicates $>$30\% missing values. High missingness is mostly structural (e.g., optional notes, category-specific fields, or particular regulatory events).
}
\end{table}

Table~\ref{tab:datasets} summarises the datasets and their most relevant attributes, including those with high missingness.

% ------------------------------------------------------------
\subsubsection*{Drugs@FDA}
% ------------------------------------------------------------
Drugs@FDA contains 51,864 product records across 29,377 applications. The product table contains only 8,388 distinct drug names and 2,878 distinct active-ingredients, showing substantial reuse across products.  The dataset is strongly imbalanced. Among applications, 78.3\% are ANDAs, 20.1\% NDAs, and 1.7\% BLAs. Three attributes exceed 30\% missingness: \texttt{ApplPublicNotes} (99.99\%), \texttt{SubmissionsPublicNotes} (99.83\%), and \texttt{ReviewPriority} (41.73\%). 

% ------------------------------------------------------------
\subsubsection*{EMA Medicines}
% ------------------------------------------------------------
The EMA dataset contains 2,746 medicine records and 1,870 distinct active ingredient, with no duplicate rows. It is predominantly composed of human medicines (85.6\%). Medicine status is also imbalanced: 68.5\% are currently authorised, 15.4\% withdrawn, and 10.4\% represent withdrawn applications. Fifty medicine names occur more than once, again showing that names alone do not uniquely identify entities. Twelve EMA attributes exceed 30\% missingness, mostly due to structural reasons (e.g., veterinary attributes being inapplicable to human medicines). 

% ------------------------------------------------------------
\subsubsection*{openFDA NDC Directory}
% ------------------------------------------------------------
The openFDA National Drug Code (NDC) Directory contains product listing information submitted to the FDA by drug labelers. It contains 138,209 product records and 7,682 distinct active ingredient entries. The listing includes both FDA-approved and unapproved products. Across the 10 selected columns, 113,976 of 1,382,090 values are missing (approx. 8.25\%). Notable missingness includes \texttt{application\_number} (27.56\%), \texttt{manufacturer} (20.30\%), and \texttt{route} (17.07\%).

\subsection{Integrated schema and overlap with input schemata}

Because the entity is the active ingredient, the integrated schema aggregates product-level information from all three sources into one record per ingredient. Required transformations include splitting the Drugs@FDA Form field and the NDC active ingredients field, breaking down combination products, and normalising NDC application numbers. The integrated schema comprises 10 attributes (Table~\ref{tab:schema}).

\begin{table}[H]
\centering
\caption{Attribute Intersection with Integrated Schema}
\label{tab:schema}
\small
\begin{tabular}{|l|l|>{\raggedright\arraybackslash}p{6cm}|}
\hline
\textbf{Attribute name} & \textbf{Attribute type} & \textbf{Datasets in which the attribute is found} \\
\hline
\texttt{activeIngredientName} & string & Drugs@FDA, EMA Medicines, openFDA NDC \\
\texttt{brandNames} & list (string) & Drugs@FDA, EMA Medicines, openFDA NDC \\
\texttt{manufacturers} & list (string) & Drugs@FDA, EMA Medicines, openFDA NDC \\
\texttt{marketingStatus} & string (categorical) & Drugs@FDA, EMA Medicines, openFDA NDC \\
\texttt{dosageForms} & list (string) & Drugs@FDA, openFDA NDC \\
\texttt{routes} & list (string) & Drugs@FDA, openFDA NDC \\
\texttt{strengths} & list (string) & Drugs@FDA, openFDA NDC \\
\texttt{applicationNumbers} & list (string) & Drugs@FDA, openFDA NDC \\
\texttt{usApprovalDate} & date & Drugs@FDA \\
\texttt{euAuthorisationDate} & date & EMA Medicines \\
\hline
\end{tabular}
\end{table}

\texttt{activeIngredientName} is the primary matching key. \texttt{brandNames} and \texttt{manufacturers} occur in all three datasets and are modelled as lists (fusion by voting). \texttt{marketingStatus} assigns each ingredient to one of five classes, derived from source-specific status flags (fused by precedence: Active $>$ Suspended $>$ Pending $>$ Discontinued $>$ Not authorised). \texttt{usApprovalDate} and \texttt{euAuthorisationDate} describe distinct regulatory events and stay separate.

% ============================================================
\subsection{Entity Count Requirement}
% ============================================================

The three datasets collectively cover a broad and overlapping part of the pharmaceutical market. Drugs@FDA and EMA Medicines focus on formally approved drugs, so substances approved in both the US and EU are represented in both sources. The openFDA NDC Directory adds coverage of US product listings, including approved and marketed products. Matching normalized \texttt{name + active substance} identifies 428 entities shared by the FDA and EMA, while matching normalized names alone identifies 752; neither result demonstrates the required minimum of 1,000 shared entities.

The openFDA NDC Directory, the larger US source, could help close this gap. Since both FDA datasets cover US products, substantial overlap between them is expected: a 30\% match rate against Drugs@FDA entries would yield more than 1,000 matches between the US sources alone. Similarity-based matching may also recover pairs missed by exact normalization, including salt-form variants and biosimilar suffixes. 

\end{document}